\chapter{体論}
この章では体論・ガロア理論についてまとめます. 代数学IIでは, 分離拡大や正規拡大など, 体拡大に関する基本事項を学び, とくにガロア拡大に現れる強い対称性を調べました. また, 有限次ガロア拡大の中間体とガロア群の部分群を対応させる, ガロア理論の基本定理を学習しました. 体は一般に無限集合であり複雑に見えますが, ガロア群の位数は拡大次数と一致するため, 有限群論の枠組みで調べられます. このため, 有限群論的な現象と体拡大に関する現象を相互に対応させることができます. \par
体拡大としては$\mb{Q}$上の拡大(代数体)に関することが多いが, もっと特殊な体拡大(たとえば有限体, 円分体など)に関することも学習した. 円分体のガロア群は$(\mb{Z}/n\mb{Z})^{\ast}$の形であり, 有限体の拡大はフロベニウス写像によって生成される巡回群である, という事実もたいへん重要である.
\par
問題演習においては, $\mb{Q}$上の多項式の最小分解体を調べ, そしてそのガロア群について調べるという問題が非常に多かった. まず拡大次数を求めることが鉄則だが, そのためには$\mb{Q}$上の多項式の既約性を求めることがよくある. そのためにEisensteinの既約判定法やガウスの補題など, 環論的ではあるがさまざまな命題について学習する必要がある.
